\documentclass[11pt]{article}
\usepackage{amsmath,amssymb,amsthm}
\usepackage[margin=1in]{geometry}
\usepackage{hyperref}
\usepackage{listings}
\usepackage{xcolor}

\lstset{
  basicstyle=\ttfamily\small,
  keywordstyle=\color{blue},
  commentstyle=\color{gray},
  showstringspaces=false,
  breaklines=true,
}

\newtheorem{theorem}{Theorem}
\newtheorem{lemma}[theorem]{Lemma}
\newtheorem{definition}[theorem]{Definition}
\newtheorem{remark}[theorem]{Remark}
\newtheorem{proposition}[theorem]{Proposition}

\title{Formalization of G\"odel's Incompleteness Theorems\\
  in Basic Recursive Arithmetic (BRA)}
\author{Summary of an Agda development}
\date{May 2026}

\begin{document}
\maketitle

\begin{abstract}
We present a fully formalized proof of both G\"odel's First and Second
Incompleteness Theorems in a tree-based variant of Guard's Basic
Recursive Arithmetic (BRA), executed in Agda~2.8 with
\verb|--safe --without-K --exact-split|. The development contains
\textbf{zero postulates and zero holes}; every per-file solo
typecheck completes in under one second (with one historical
exception, the 16k-line dispatcher cascade for $\mathit{thmT}$).

The headline results are:
\begin{itemize}
  \item $\mathit{thm11} : \mathit{Deriv}\;G \to \mathit{Deriv}\;\bot$
        (G\"odel I: provability of the diagonal sentence implies $\bot$);
  \item $\mathit{godelII} : \mathit{Deriv}\;\mathit{Con} \to \mathit{Deriv}\;\bot$
        (G\"odel II: provability of consistency implies $\bot$).
\end{itemize}
The development follows Guard's 1963 lecture notes on recursive
arithmetic exclusively.  The intermediate construction is a single
first-order computable function $\mathit{constr5} : \mathit{Fun1}$
realizing Guard's ``$F(x)$'' (page~17 of his notes); each step in the
chain is given by an explicit equational derivation in BRA.
\end{abstract}

\section{Introduction}

The aim of the project is to give a feasible, ASCII-only, syntax-first
Agda formalization of G\"odel's two incompleteness theorems, following
Guard's 1963 system of \emph{Basic Recursive
Arithmetic}~\cite{guard1963}.  The intent is to avoid historical
baggage (G\"odel numbering of Hilbert proofs, primitive recursive
arithmetic with explicit numerals, ``encoded'' substitution as a
Hilbert-system meta theorem) by working inside a small tree-based
equational system from the outset:

\begin{itemize}
  \item Terms are built from $O$ (leaf), $\mathit{var}\;n$, $\mathit{ap}_1\;f\;t$, $\mathit{ap}_2\;g\;t_1\;t_2$.
  \item Function symbols come in two arities: $\mathit{Fun}_1$ and $\mathit{Fun}_2$, with primitive constants
        $\{I, \mathit{Fst}, \mathit{Snd}, \mathit{Comp}, \mathit{Comp}_2, \mathit{Z}, \mathit{KT}\,\_\}$ and
        $\{\mathit{Pair}, \mathit{Const}, \mathit{Lift}, \mathit{Post}, \mathit{Fan}, \mathit{IfLf}, \mathit{TreeEq}, \mathit{treeRec}\}$.
  \item Equations are objects $\mathit{eqn}\;\ell\;r$ between terms; formulas $\mathit{Formula}$ extend equations with $\mathit{not}$ and $\mathit{imp}$.
  \item Provability is an inductive type $\mathit{Deriv}\;F$ with constructors for the algebraic axioms ($\mathit{ax}_*$), the equality axioms,
        the propositional axioms ($\mathit{ax}_K, \mathit{ax}_S, \mathit{ax}_{\mathit{Neg}}, \mathit{ax}_{\mathit{ExFalso}}, \mathit{ax}_{\mathit{Contrapos}}$), modus ponens ($\mathit{mp}$),
        substitution ($\mathit{ruleInst}$), and binary tree induction ($\mathit{ruleIndBT}$).
\end{itemize}

The natural numbers are represented as binary trees ($\mathit{lf}$ for $0$, $\mathit{nd}\;\mathit{lf}\;n$ for the successor in unary).
The encoding $\mathit{reify} : \mathit{Tree} \to \mathit{Term}$ maps $\mathit{lf} \mapsto O$ and $\mathit{nd}\;a\;b \mapsto \mathit{ap}_2\;\mathit{Pair}\;(\mathit{reify}\;a)\;(\mathit{reify}\;b)$.

\section{The provability predicate $\mathit{thmT}$}

The G\"odel-style ``$\mathit{th}$'' provability predicate is realized in BRA by a single $\mathit{Fun}_1$-functor:
\[
  \mathit{thmT} \;=\; \mathit{Rec}\;\mathit{stepProtoWrapped} \;:\; \mathit{Fun}_1.
\]
Here $\mathit{stepProtoWrapped}$ is a $\mathit{Fun}_2$ assembled from a 40-deep linear $\mathit{IfLf}$ cascade
(\verb|BRA/Thm/StepProto.agda| and \verb|BRA/Thm/ThmT.agda|, $\sim$16k lines), one branch per primitive proof rule:
ten equational axioms of Group~I (\verb|axI|, \verb|axFst|, \dots, \verb|axKT|), eight Group~II axioms
(\verb|axIfLfL|, \verb|axIfLfN|, four \verb|axTreeEq*|, \verb|axGoodstein|), nine equality and propositional rules,
and three rules of inference ($\mathit{mp}$, $\mathit{ruleInst}$, $\mathit{ruleIndBT}$).

The \emph{key} property of $\mathit{thmT}$ is that on the encoded form of any proof tree it returns the
encoded conclusion of that proof.  This is the content of the \emph{encoding-soundness} lemma
$\mathit{encode} : (P : \mathit{Formula}) \to \mathit{Deriv}\;P \to \Sigma\;\mathit{Tree}\;(\lambda y \to \mathit{Deriv}\;(\mathit{thmT}\;(\mathit{reify}\;y) = \mathit{reify}\;(\mathit{codeFormula}\;P)))$,
proved by induction on the $\mathit{Deriv}$ derivation in
\verb|BRA/Thm/ThmT.agda|.

\subsection{Self-substitution closure}

A second key property is that $\mathit{thmT}$ contains no free variable (it is a closed combinator):
\[
  \mathit{thClosed} : \forall x\;r,\quad \mathit{substF}_1\;x\;r\;\mathit{thmT} = \mathit{thmT}
\]
and that the encoding of $\mathit{thmT}$ is var-free:
\[
  \mathit{codeF1Th\_noVar} : \forall x\;\mathit{repl},\quad \mathit{codedSubst}\;\mathit{repl}\;(\mathit{code}\;(\mathit{var}\;x))\;(\mathit{codeF}_1\;\mathit{thmT}) = \mathit{codeF}_1\;\mathit{thmT}.
\]
Both are \verb|refl| inside the abstract block where $\mathit{thmT}$ unfolds.

\section{G\"odel's First Incompleteness Theorem (Theorem~11)}

\subsection{The diagonal sentence}

Following Guard, define
\begin{align*}
  F_{\mathrm{pre}} \;&:=\; \neg \big(\mathit{thmT}\;(v_1) \;=\; \mathit{sub}\;(v_0)\;(v_0)\big),\\
  j \;&:=\; \mathit{codeFormula}\;F_{\mathrm{pre}} \quad \text{(a closed Tree)},\\
  G \;&:=\; F_{\mathrm{pre}}[v_0 := \mathit{reify}\;j] \;=\; \neg\big(\mathit{thmT}\;(v_1) \;=\; \mathit{sub}\;i\;i\big),
\end{align*}
where $i := \mathit{reify}\;j$ and $\mathit{sub}$ is the BRA-internal coded substitution function (\verb|BRA/Sub.agda|).
The sentence $G$ has $v_1$ free; informally, $G$ asserts of itself that no proof tree $v_1$ proves the formula whose code is $j[v_0 := i]$.

\subsection{Theorem 11}

Pulling apart the abstract closure (\verb|BRA/Thm11Abstract.agda|), the precise statement is:

\begin{theorem}[\texttt{thm11}, G\"odel I]
\label{thm:godel-I}
\(
  \mathit{thm11} : \mathit{Deriv}\;G \to \mathit{Deriv}\;\bot,
\)
where $\bot := \mathit{atomic}\;(\mathit{eqn}\;\mathit{trueT}\;\mathit{falseT})$ with $\mathit{trueT} := O$ and $\mathit{falseT} := \mathit{ap}_2\;\mathit{Pair}\;O\;O$.
\end{theorem}

\begin{proof}[Sketch]
From any $\pi : \mathit{Deriv}\;G$ extract its encoding $(y_d, E_1)$ where
$E_1 : \mathit{Deriv}\;(\mathit{thmT}\;(\mathit{reify}\;y_d) = \mathit{reify}\;(\mathit{codeFormula}\;G))$.
A separate \emph{diagonal lemma} \texttt{diagonal} (\verb|BRA/Thm11Diagonal.agda|)
yields $N : \mathit{Deriv}\;\neg(\mathit{thmT}\;(\mathit{reify}\;y_d) = \mathit{reify}\;(\mathit{codeFormula}\;G))$
from $\pi$, and $\mathit{ax}_{\mathit{ExFalso}}$ + two $\mathit{mp}$'s collapse $E_1$ and $N$ to $\mathit{Deriv}\;\bot$.
\end{proof}

The diagonal lemma is constructed inside an abstract block in
\verb|BRA/Thm11Diagonal.agda| using $\mathit{thClosed}$, $\mathit{codeF1Th\_noVar}$,
and $\mathit{cc}$ (the local instance of the meta-equation $\mathit{codedSubst}\;(\mathit{code}\;i)\;(\mathit{code}\;v_0)\;j = \mathit{codeFormula}\;G$,
which is provable by induction on $F_{\mathrm{pre}}$ via $\mathit{codeCommutesFormula}$ in \verb|BRA/CodeCommutes.agda|).

\section{G\"odel's Second Incompleteness Theorem (Theorem~14)}

The second theorem is the more substantial construction.  It says that
the consistency formula $\mathit{Con}$ entails $\bot$:

\begin{theorem}[\texttt{godelII}, G\"odel II]
\label{thm:godel-II}
\(
  \mathit{godelII} : \mathit{Deriv}\;\mathit{Con} \to \mathit{Deriv}\;\bot,
\)
where $\mathit{Con} := \neg\big(\mathit{thmT}\;(v_0) = \mathit{reify}\;(\mathit{codeFormula}\;\bot)\big)$
(consistency: ``no proof, encoded as $v_0$, derives $\bot$'').
\end{theorem}

The construction follows Guard's pages~16--17 step by step.

\subsection{Caveat: only the \emph{coded} form is a well-defined BRA object}
\label{subsec:underline-caveat}

A structurally crucial point, easy to miss in Guard's notation, is that
the expression
\[
  f(\underline{x}) \;=\; \underline{f(x)}
\]
is \textbf{not} a BRA formula.  Guard's underline $\underline{(\cdot)}$
takes a value and returns its numeral term, an operation that is
well-defined only at the meta-level for closed numerical inputs.  For
a variable $x$ --- which is precisely the case Theorems~12--14 actually
need --- neither $\underline{x}$ nor $\underline{f(x)}$ has any BRA
referent: there is no operator in BRA's term language that, given an
arbitrary $\mathrm{Term}$, returns ``the numeral of its value.''  No
attempt at writing the equation directly as
$\mathit{atomic}\;(\mathit{eqn}\;\ldots\;\ldots)$ for a variable $x$
succeeds, because the slots cannot be filled with BRA terms.

The only formal object that \emph{does} correspond to Guard's
expression is its \emph{code} --- the Tree
\[
  \ulcorner f(\underline{x}) \;=\; \underline{f(x)} \urcorner,
\]
which Definition~12 of \texttt{guard15.pdf} \emph{defines} by
stipulation.  In our Agda development this Tree is built parametrically
as $\mathit{codeFTeq}_1\;f\;x$, with $\mathit{cor}\;x$ filling the slot
for the code of $x$'s numeral and $\mathit{cor}\;(f\,x)$ filling the
slot for the code of $f(x)$'s numeral.

\textbf{Consequence for Theorem~12's statement.}  When we write
\[
  \mathit{Deriv}\,\big(\mathit{thmT}\;(D_f\;x)
       \;=\; \mathit{codeFTeq}_1\;f\;x\big),
\]
the right-hand side is a \emph{specific Tree}, not a BRA formula.  The
colloquial reading ``BRA-internally, $f(\underline{x}) =
\underline{f(x)}$'' is shorthand for ``BRA proves that the proof-checker
$\mathit{thmT}$ applied to $D_f\;x$ equals the Tree named by
Definition~12.''  Throughout the rest of this section, every occurrence
of $f(\underline{x}) = \underline{f(x)}$ (and the analogous
$f(\underline{x}) = \underline{y}$ form used in Theorem~13) should be
read as a name for the defined Tree, never as a free-standing BRA
formula.

\subsection{Numerals: the asymmetry $f(\underline{x})$ vs.\ $\underline{f(x)}$}
\label{subsec:num-cor}

A crucial feature of the Theorem~12--14 chain, easy to miss in the
formalisation, is the \emph{asymmetric} role of the numeral function.
In Guard's notation~\cite{guard1963} the underline $\underline{x}$
denotes the \emph{numeral} of $x$: a closed BRA term whose value is
$x$.  In our Agda development this is the primitive-recursive function
\[
  \mathit{num} \;=\; \mathit{cor} \;:\; \mathit{Tree} \to \mathit{Tree},
\]
which sends a tree $x$ to the code of the BRA term that \emph{names}
$x$.  We use $\mathit{num}$ and $\mathit{cor}$ interchangeably with
Guard's $\underline{x}$.

The asymmetry is this: for any $\mathit{Fun}_1$-functor $f$ and tree $x$,
\begin{itemize}
  \item $\mathit{ap}_1\;f\;(\mathit{num}\;x)$, i.e.\ $f(\underline{x})$,
        \emph{is} the code of a BRA term --- specifically the term
        ``$f$ applied to the numeral of $x$'';
  \item $\mathit{num}\;(f\;x)$, i.e.\ $\underline{f(x)}$, \emph{is generally not}
        the code of any such syntactic term --- it is the code of the numeral
        naming the \emph{value} $f\;x$.
\end{itemize}
The two trees disagree in general; they agree only after BRA proves
them equal.

Theorem~12 is precisely the internalisation of that equation: for
every $f : \mathit{Fun}_1$ it builds a $D_f : \mathit{Fun}_1$ such that
\[
  \mathit{Deriv}\;\big(\mathit{thmT}\;(D_f\;x)
       \;=\; \mathit{codeFTeq}_1\;f\;x\big),
\]
where the right-hand side $\mathit{codeFTeq}_1\;f\;x$ is the Tree
defined by Guard's Definition~12 as the code of
$f(\underline{x}) = \underline{f(x)}$ --- recall from
Section~\ref{subsec:underline-caveat} that this Tree, not the
unquoted equation, is the formal object.  Its outer shape is a
$\mathit{Pair}$ of the code-of-LHS (encoding ``$\mathit{ap}_1\;f$
applied to the code of $x$'s numeral'', with $\mathit{cor}\;x$ in
the inner slot) and the code-of-RHS (the code of $f(x)$'s numeral,
which is $\mathit{cor}\;(f\,x)$).

In other words: BRA itself proves that the proof-checker
$\mathit{thmT}$, applied to $D_f\;x$, equals the Tree named by
Guard's Definition~12.  This is what licenses every subsequent step
of Goedel~II: it lets the internal $\mathit{thmT}$ predicate move
between the syntactic and semantic sides of the encoded equation
exactly where Guard's pencil-and-paper argument tacitly does so.

A symmetric variant in which both sides are coded the same way
($\mathit{code}\;(f\;x)$ on both sides, say) collapses: the resulting
trees do not align across the five steps of Theorem~14, and step~5
becomes underivable.  The asymmetric form $f(\underline{x}) =
\underline{f(x)}$ is therefore not a presentational choice but a load-bearing
feature of the proof.

\subsection{Reading the slot \texttt{cor u} in encoded shapes}
\label{subsec:cor-slot-reading}

A second subtlety hidden by Guard's underline notation concerns the
slot occupied by \texttt{cor u}.  Pure \verb|code| of a BRA term is
defined by
\begin{verbatim}
  code (ap1 f u) = nd tagAp1 (nd (codeF1 f) (code u)),
\end{verbatim}
whereas Guard's $\widehat{f(\underline{u})}$ uses \verb|cor| in the
inner slot:
\begin{verbatim}
  code "f u_"   = nd tagAp1 (nd (codeF1 f) (cor u)).
\end{verbatim}
The notation tempts the reader to think this Tree is itself
\verb|code t| for some BRA term \verb|t|, but in general it is not.

\emph{When the slot is filled by a syntactic \texttt{ap1 cor x} for a
free variable \texttt{x}.}  This is the form that actually appears in
the Theorem~14 chain (e.g.\ \verb|encoded_th_x_at| in
\verb|BRA/Th14Step3.agda|, with \verb|x : Term| held abstract,
typically \verb|x = var 0|).  Read syntactically, the assembled
expression IS \verb|code t| for
\[
  t \;=\; \mathit{ap}_1\,f\,(\mathit{ap}_1\,\mathit{cor}\,x):
\]
the BRA term ``$f$ applied to $\mathit{cor}\,x$'' --- $\mathit{cor}$
kept syntactic, the variable position passed through unevaluated.
This is the only sense in which $\widehat{f(\underline{u})}$ at a
\emph{variable} $u$ literally encodes a single BRA term.

\emph{An example where the slot expression is not \texttt{code t} for
any \texttt{t}.}  Trees of the same outer shape need not be codes of
BRA terms.  Take e.g.\
\begin{verbatim}
  nd tagAp1 (nd (codeF1 f) (nd tagAp1 lf)).
\end{verbatim}
For \verb|nd tagAp1 lf| to be \verb|code t'|, we would need
\verb|t' = ap1 g s| with \verb|nd (codeF1 g) (code s) = lf|, which is
impossible (the payload is never \verb|lf|).  So the inner slot
\verb|nd tagAp1 lf| is not the code of any BRA term, and consequently
the whole Tree above is not \verb|code| of anything.  Non-code Trees
of this kind do arise as intermediate values inside the \verb|thmT|
dispatcher cascade --- the encoding-shape is not protected against
them by the type alone.

\emph{Bridging the two readings.}  At a closed numeric Tree \verb|u|,
\verb|corOfReify| (\verb|BRA/Cor.agda|) gives
\verb|cor u = code (reify u)| internally as a \texttt{Deriv}, so under
that reading the assembled Tree equals \verb|code (ap1 f (reify u))|
--- a \emph{different} \verb|t| (the term ``$f$ applied to the numeral
of $u$'') from the variable reading above.  Theorem 12 is precisely
the BRA-internal lemma that equates the two: the equation
$\mathit{ap}_1\,f\,(\mathit{cor}\,x) = \mathit{cor}\,(f\,x)$ is the
internal version of ``the syntactic-variable reading and the
evaluated-numeral reading agree''.

\subsection{Theorem 12 --- the encoded ``$\mathit{th}\;f$'' clause}

For every $\mathit{Fun}_1$-functor $f$, Theorem~12 produces a $\mathit{Fun}_1$
$D_f$ such that under the canonical encoding,
$\mathit{thmT}$ evaluated at $D_f\;x$ is provably equal to the encoded form of
``$f(\mathit{num}\;x) = \mathit{num}\;(f\;x)$''.  Formally:

\begin{theorem}[$\mathit{thm12}\_F_1$, parametric Theorem 12 for $\mathit{Fun}_1$]
\(
  \mathit{thm12}\_F_1 : (f : \mathit{Fun}_1) \to \mathit{Thm12\_F_1\_Result}\;f
\)
where the record $\mathit{Thm12\_F_1\_Result}\;f$ packages a $D_f : \mathit{Fun}_1$ together with
\(
  \mathit{proof} : (x : \mathit{Term}) \to \mathit{Deriv}\;(\mathit{thmT}\;(D_f\;x) = \mathit{codeFTeq}_1\;f\;x).
\)
The conclusion $\mathit{codeFTeq}_1\;f\;x$ is the literal Pair-shape encoded equation
$\mathit{ap}_1\;f\;(\mathit{cor}\;x) = \mathit{cor}\;(f\;x)$.
\end{theorem}

(Analogous statement for $\mathit{Fun}_2$.)
This is a structural recursion on $f$: there is one clause per
$\mathit{Fun}_1$ / $\mathit{Fun}_2$ constructor (15 total), each in its own
\verb|BRA/Thm12/Parts/*.agda| file.  The composite cases ($\mathit{Comp}$,
$\mathit{Comp}_2$, $\mathit{Lift}$, $\mathit{Post}$, $\mathit{Fan}$, $\mathit{treeRec}$)
dispatch to per-case modules; the recursive $\mathit{TreeEq}$ case uses
\verb|ruleIndBT2| (2-tree-pair induction with diagonal cross-IHs) to
establish closure.  The whole theorem is in
\verb|BRA/Thm12.agda|, module \verb|FromBridges|, and contains zero
postulates / zero holes.

\subsection{Theorem 13 --- under hypothesis ``$f\;x = y$''}

Building on Theorem~12, one produces under the hypothesis $f\;x = y$
the encoded form of ``$f(\mathit{num}\;x) = \mathit{num}\;y$'':

\begin{lemma}[$\mathit{thm13}\_F_1$]
For every $f : \mathit{Fun}_1$ and every Theorem~12 result $r_{12}$ for $f$,
\(
  \mathit{thm13}\_F_1\;f\;r_{12}\;x\;y : \mathit{Deriv}\;(f\;x = y) \to
    \mathit{Deriv}\;(\mathit{thmT}\;(D_f\;x) = \mathit{codeFTeq}_{1,\mathrm{Hyp}}\;f\;x\;y)
\)
where the conclusion replaces the ``$\mathit{cor}\;(f\;x)$'' slot of
$\mathit{codeFTeq}_1$ by $\mathit{cor}\;y$.  One line: $\mathit{ruleTrans}$ on
$r_{12}.\mathit{proof}\;x$ with $\mathit{cong}_R$~$\mathit{Pair}$ on $\mathit{cong}_1\;\mathit{cor}$ of the hypothesis.
\end{lemma}

\subsection{The function \texttt{sb} beyond its specification}
\label{subsec:sb-nonstandard}

Guard (\cite{guard1963}, p.~14, Exercise~24~[4]) specifies
\verb|sb| as a substitution function only on encoded arguments:
\[
  \mathit{sb}\big(J(\ulcorner T\urcorner, n),\, \ulcorner P\urcorner\big)
    \;=\; \ulcorner P[T/x_n]\urcorner.
\]
That is, the lecture notes describe \verb|sb| only when the first slot
is the code $\ulcorner T\urcorner$ of a BRA term \verb|T|.  In our
formalisation \verb|sb : Fun2| is defined uniformly via the underlying
tree-level substitutor \verb|codedSubst| of \verb|BRA/Term.agda|:
\begin{verbatim}
  codedSubst : Tree -> Tree -> Tree -> Tree
  codedSubst repl varCode lf       = lf
  codedSubst repl varCode (nd a b) = boolCase (treeEq varCode (nd a b))
                                       repl
                                       (nd (codedSubst repl varCode a)
                                           (codedSubst repl varCode b))
\end{verbatim}
The substituent \verb|repl| here is just \emph{any} Tree; nothing in
\verb|codedSubst| relies on \verb|repl = code T| for some BRA term
\verb|T|.  The corresponding spec (\verb|BRA/Sb.agda|, \verb|sb_def|)
\begin{verbatim}
  ap2 sb (ap2 Pair (reify codeA) (reify (natCode n))) (reify codeP)
    = reify (codedSubst codeA (code (var n)) codeP)
\end{verbatim}
holds for any Tree \verb|codeA|, not just for \verb|codeA = code T|.
Guard's lecture notes do not state this generalisation.

\emph{Crucial use in Theorems 13--14.}  In \verb|BRA/Th14Step4.agda|
the substituent slot is filled by \verb|ap1 cor x| with
\verb|x : Term| held abstract:
\begin{verbatim}
  ap2 sb (ap2 Pair (ap1 cor x) varCode1T) (reify (codeFormula G)).
\end{verbatim}
By Section~\ref{subsec:cor-slot-reading} the value \verb|cor x| at a
closed numeric \verb|x| is \emph{not} \verb|code T| for any BRA term
\verb|T|: it is a numeral encoding produced by \verb|cor|'s recursion
on the binary-tree shape.  Yet \verb|sb| processes it cleanly,
walking \verb|codeFormula G| and replacing each subtree equal to
\verb|code (var 1)| by the raw Tree \verb|cor x|.  The whole
construction of $K\_\mathit{part}$, $g\_\mathit{part}$, and
$\mathit{constr5}$ in Theorem~14 (Section~\ref{thm:godel-II} below)
depends essentially on this --- the Definition~12 line~2 of $\mathit{th}$
(\verb|guard15.pdf| p.\ 15) substitutes a \emph{numeral} into a coded
formula, and a numeral is not the code of any BRA term.

\subsection{Theorem 14 --- the consistency-to-bot construction}

This is the heart of G\"odel II. The contract abstracted in
\verb|BRA/Thm14Abstract.agda| asks for:
\begin{enumerate}
  \item a $\mathit{Fun}_1$-functor $\mathit{constr5}$ such that $\mathit{thmT}\;(\mathit{constr5}\;x)$ syntactically encodes $\bot$;
  \item a parametric internal-implication
\[
  \mathit{step5} : (x : \mathit{Term}) \to
   \mathit{Deriv}\big((\mathit{thmT}\;x = \mathit{reify}\;(\mathit{codeFormula}\;G)) \;\mathit{imp}\; (\mathit{thmT}\;(\mathit{constr5}\;x) = \mathit{reify}\;(\mathit{codeFormula}\;\bot))\big);
\]
\end{enumerate}

The closure $\mathit{closureToG}$ then applies $\mathit{ax}_{\mathit{Contrapos}}$ to $\mathit{step5}$ instantiated at $v_1$,
$\mathit{ruleInst}$s $\mathit{Con}$ at $\mathit{constr5}\;v_1$, modus-ponenses to obtain $\neg(\mathit{thmT}\;v_1 = i)$,
transports through $\mathit{subIIeqcG}$ ($\mathit{sub}\;i\;i = \mathit{cor}\;G$) and $G_{\mathrm{unfold}}$
(propositional normal form of $G$), arriving at $\mathit{Deriv}\;G$.  Composing with $\mathit{thm11}$ gives
$\mathit{godelII} : \mathit{Deriv}\;\mathit{Con} \to \mathit{Deriv}\;\bot$.

\subsubsection{Construction of $\mathit{constr5}$}

Following Guard's page~17, the construction relies on two
pre-existing tautologies, encoded:
\begin{itemize}
  \item $t := (v_0 = v_2)\;\mathit{imp}\;((v_1 = v_2)\;\mathit{imp}\;(v_0 = v_1))$ (transitivity);
  \item $t' := (v_0 = v_1)\;\mathit{imp}\;((\neg(v_0 = v_1))\;\mathit{imp}\;\bot)$ (ex falso).
\end{itemize}
Both are derivable in BRA (\verb|BRA/Thm14Numerals.agda|), and via $\mathit{encode}$ they yield
closed Trees $\mathit{tterm}, \mathit{t'term}$ such that $\mathit{thmT}\;(\mathit{reify}\;\mathit{tterm}) = \mathit{reify}\;(\mathit{codeFormula}\;t)$
(and analogously for $t'$).

The chain is:
\begin{align*}
  f\_\mathit{part}(x) \;&:=\; \mathit{ruleInst}\;v_0\;(\mathit{th}\;\hat x)\;\big(\mathit{ruleInst}\;v_1\;(\mathit{sub}\;i\;i)\;(\mathit{ruleInst}\;v_2\;\mathit{cor}\;G\;\mathit{tterm})\big),\\
  g\_\mathit{part}(x) \;&:=\; \mathit{mp}\big(\mathit{mp}(f\_\mathit{part}(x), D_{\mathit{thmT}}\;x), D_{\mathit{sub}}\;i\;i\big),\\
  K\_\mathit{part}(x) \;&:=\; \mathit{ruleInst}\;v_1\;(\mathit{cor}\;x)\;x,\\
  h\_\mathit{part}\_\mathit{skel}(x) \;&:=\; \mathit{ruleInst}\;v_0\;(\mathit{th}\;\hat x)\;\big(\mathit{ruleInst}\;v_1\;(\mathit{sub}\;i\;i)\;\mathit{t'term}\big),\\
  \mathit{constr5}_{\mathrm{final}}(x) \;&:=\; \mathit{mp}\big(\mathit{mp}(h\_\mathit{part}\_\mathit{skel}(x), g\_\mathit{part}(x)), K\_\mathit{part}(x)\big).
\end{align*}
The substituents are the encoded mixed-form Pairs of Guard:
$\widehat{\mathit{th}\,\hat x} := \mathit{ap}_2\;\mathit{Pair}\;(\mathit{reify}\;\mathit{tagAp}_1)\;(\mathit{ap}_2\;\mathit{Pair}\;(\mathit{reify}\;(\mathit{codeF}_1\;\mathit{thmT}))\;(\mathit{cor}\;x))$,
$\mathit{sub}_i\widehat{\,}_i := \mathit{reify}\;(\mathit{code}\;(\mathit{ap}_2\;\mathit{sub}\;i\;i))$, and $\mathit{cor}\;G := \mathit{reify}\;(\mathit{code}\;(\mathit{reify}\;(\mathit{codeFormula}\;G)))$.

\subsubsection{The five steps and their lemmas}

Following Guard's enumeration (\verb|guard15.pdf| pages~16--17):

\begin{description}
  \item[Step 1 --- $\mathit{step1\_l}$] (\verb|Th14Constr5.agda|).  By Theorem~13 applied to $\mathit{thmT}$ at $(x, \mathit{cor}\;G)$ under hypothesis $\mathit{thmT}\;x = \mathit{cor}\;G$, get
  $\mathit{Deriv}\;(\mathit{PrAtX}\;x \mathrel{\mathit{imp}} \mathit{thmT}\;(D_{\mathit{thmT}}\;x) = \mathit{codeFTeq}_{1,\mathrm{Hyp}}\;\mathit{thmT}\;x\;\mathit{cor}\;G)$.

  \item[Step 2 --- $\mathit{step2\_l}$] (\verb|Th14Step2.agda|).  By Theorem~13 applied to $\mathit{sub}$ at $(i, i, \mathit{cor}\;G)$ under closed hypothesis $\mathit{subIIeqcG}$.

  \item[Step 3 --- $\mathit{step3\_l}$ (= $g\_\mathit{part}\_l$)] (\verb|Th14Step3.agda|, \verb|Th14Step3Canon.agda|).  Three sb-applications on $t$ via $\mathit{thmTSubLemma}$, canonicalized via $\mathit{subT}\_\mathit{codeFormula}\_\mathit{imp/atomic}$ and leaf preservation lemmas, then two $\mathit{mp}$ dispatches at $g\_\mathit{part}\_\mathit{inner}$ and $g\_\mathit{part}$.  Output: $\mathit{Deriv}\;(\mathit{PrAtX}\;x \mathrel{\mathit{imp}} \mathit{thmT}\;(g\_\mathit{part}\;x) = \widehat{\mathit{th}\;\hat x = \mathit{sub}\;i\;i})$.

  \item[Step 4 --- $K\_\mathit{part}\_l$] (\verb|Th14Step4Final.agda|).  One sb-application at $v_1$ on the diagonal sentence $G_{\mathrm{norm}}$ via $\mathit{thmTSubLemma}$ + the $\mathit{PrAtX}\;x$ hypothesis to swap $\mathit{thmT}\;x \mapsto \mathit{cor}\;G$, canonicalized via $\mathit{subT}\_\mathit{codeFormula}\_\mathit{neg/atomic}$ and leaf preservation.  Output: $\mathit{Deriv}\;(\mathit{PrAtX}\;x \mathrel{\mathit{imp}} \mathit{thmT}\;(K\_\mathit{part}\;x) = \widehat{\mathit{th}\;\hat x \neq \mathit{sub}\;i\;i})$.

  \item[Step 5 --- $\mathit{step5\_l}$] (\verb|Th14Step5.agda|).  Two outer $\mathit{mp}$ dispatches at $\mathit{constr5}\_{\mathrm{inner,\,final}}$ and $\mathit{constr5}_{\mathrm{final}}$, with $\mathit{step3\_l}$ feeding the inner antecedent and $K\_\mathit{part}\_l$ standing in syntactically as the second-mp witness.  The dispatcher's syntactic projection of the consequent gives the encoded $\bot$.  Output: $\mathit{Deriv}\;(\mathit{PrAtX}\;x \mathrel{\mathit{imp}} \mathit{thmT}\;(\mathit{constr5}_{\mathrm{final}}\;x) = \mathit{reify}\;(\mathit{codeFormula}\;\bot))$.

  \item[Auxiliary --- $h\_\mathit{part}\_\mathit{skel}\_l$] (\verb|Th14Step5HPart.agda|).  Two sb-applications on $t'$, mirror of step~3 (one fewer layer because $t'$ has only two free variables).  Output: $\mathit{thmT}\;(h\_\mathit{part}\_\mathit{skel}\;x) = \widehat{\text{ex falso form}}$.  Used as $\mathit{step5\_l}$'s inner $d_1$.
\end{description}

The headline composition (\verb|BRA/GoedelIIFull.agda|) is then:

\begin{verbatim}
r12_th  : Thm12_F1_Result thmT
r12_th  = FromBridges.thm12_F1 thmT

r12_sub : Thm12_F2_Result sub
r12_sub = FromBridges.thm12_F2 sub

module Step5 = BRA.Th14Step5.Th14Step5 r12_th r12_sub
module Final = BRA.GoedelII.Compose Step5.constr5_final Step5.step5_l

godelII : Deriv Con -> Deriv bot
godelII = Final.godelII
\end{verbatim}

\subsection{What Theorems 12--14 are doing, and where the notation hides things}
\label{subsec:modal-perspective}

The BRA chain involves two distinct encodings of trees as trees.  Pulling
them apart helps explain what each of Theorems 12--14 is actually doing,
and why Guard's compact notation can mislead.

\paragraph{Two encodings.}
Two functions are at play, and they are not the same:
\begin{itemize}
  \item $\mathit{cor} : \mathrm{Tree} \to \mathrm{Tree}$, the
        \emph{numeral encoding}.  $\mathit{cor}\,x$ (also written
        $\underline{x}$) is the canonical numeral whose value is $x$,
        a tree of the form $S(S(\ldots S(0)))$.
  \item $\mathit{code} : \mathrm{Term} \to \mathrm{Tree}$, the
        \emph{syntactic encoding} of a BRA term, recording the term's
        syntax tree.
\end{itemize}
The two operations agree on no input shape: $\mathit{cor}\,x$ collapses
a tree to its numerical magnitude, while $\mathit{code}\,t$ preserves
the structure of $t$ as a syntax tree.

\paragraph{The notational gotcha: $\mathit{code}(f(\underline{x}))$ vs $\mathit{code}(\underline{f(x)})$.}
Two expressions that look symmetric in Guard's compact notation are
quite different syntactic objects:
\begin{itemize}
  \item $f(\underline{x})$ is the \emph{syntactic term} ``function
        symbol $f$ applied to the numeral $\underline{x}$''.  Its
        encoding $\mathit{code}(f(\underline{x}))$ is unambiguously
        the code of a (closed) BRA term: $f$ applied to the numeral
        of $x$.
  \item $\underline{f(x)}$ is the \emph{numeral} whose value is $f(x)$
        --- a tree $S^n(0)$ for some $n$.  Guard's notation
        $\mathit{code}(\underline{f(x)})$ takes the encoding of
        \emph{that numeral}, which is in general \emph{not} the code
        of the term $f(x)$: the term $f(x)$ has shape ``$f$ applied
        to $x$'', while the numeral $\underline{f(x)}$ has shape
        $S^n(0)$.  Completely different syntactic trees.
\end{itemize}
The compact $\underline{(\cdot)}$ invites the reader to treat the two
sides of an equation like $\widehat{f(\underline{x}) = \underline{f(x)}}$
as ``the same expression with the bar moved''.  In fact the inner
brackets quietly switch between two different encodings: a
syntactic-term encoding on the left, a numeral encoding on the right.
Theorem 12 is the BRA-internal proof that, despite the notational
mismatch, these two trees evaluate to the same value at every $f$.

\paragraph{Theorem 12: a per-functor lift $D_f$ together with naturality.}
For every $\mathit{Fun}_1$ functor $f$, Theorem 12 produces a primitive
recursive $D_f : \mathit{Fun}_1$ together with an internal derivation
\[
  \mathit{thmT}\,(D_f\,x) \;=\; \widehat{f(\underline{x}) \;=\; \underline{f(x)}}.
\]
$D_f$ is the ``code-level cousin'' of $f$: it takes numerals to
numerals.  The encoded equation says that $D_f$ commutes with the
numeral encoding, in the sense that running $f$ on a numeral and taking
the numeral of $f(x)$ produce the same value.  Theorem 12 simultaneously
\emph{defines} the lift $D_f$ and \emph{proves} that the lift is faithful.

\paragraph{Theorem 13: lifting an external equation to an internal one.}
Given a value-level $\mathit{Deriv}\,(f\,x = y)$, Theorem 13 produces
an internal derivation
$\mathit{thmT}\,(D_f\,x) = \widehat{f(\underline{x}) = \underline{y}}$.
Read externally: ``BRA proves $f(x) = y$ at the value level $\Rightarrow$
BRA also internally proves the corresponding encoded equation between
numerals''.  This is the Hilbert--Bernays necessitation rule D1
($\vdash A \Rightarrow \vdash \mathrm{Pr}(\ulcorner A \urcorner)$)
restricted to equational $A$, applied on top of Theorem 12 to rewrite
the $\underline{f(x)}$ side of the encoded equation as $\underline{y}$.

\paragraph{Theorem 14, step 5: internalised G\"odel~I at the diagonal $G$.}
$\mathit{step5\_l}$ derives, parametrically in $x$, the encoded internal
implication
\[
  \mathit{thmT}\,x = \mathit{cor}\,G \;\;\to\;\;
  \mathit{thmT}\,(\mathit{constr5}\,x) = \mathit{cor}\,\bot.
\]
Read externally: ``if $x$ codes a proof of $G$, then
$\mathit{constr5}\,x$ codes a proof of $\bot$'', i.e.\ the encoded
implication $\mathrm{Pr}(\ulcorner G\urcorner) \to
\mathrm{Pr}(\ulcorner \bot\urcorner)$.

The textbook decomposition of this internal implication has three
pieces, all combined by internal modus ponens:
\begin{enumerate}
  \item One direction of the diagonal,
        $G \to \neg\,\mathrm{Pr}(\ulcorner G\urcorner)$, BRA-derivable
        from the construction of $G$.  Necessitating this gives an
        internal derivation of
        $\mathrm{Pr}(\ulcorner G\urcorner) \to
         \mathrm{Pr}(\ulcorner \neg\,\mathrm{Pr}(\ulcorner G\urcorner)\urcorner)$.
  \item Hilbert--Bernays clause D2 at $A := G$:
        $\mathrm{Pr}(\ulcorner G\urcorner) \to
         \mathrm{Pr}(\ulcorner \mathrm{Pr}(\ulcorner G\urcorner)\urcorner)$.
        In BRA the realiser of D2 at $G$ is just the proof-checker:
        given a tree $x$ with $\mathit{thmT}\,x = \mathit{cor}\,G$,
        the verification ``$x$ codes a proof of $G$'' is itself the
        evidence that $\mathrm{Pr}(\ulcorner G\urcorner)$ holds.
  \item Internal modus ponens combining the two: the conclusion
        $\mathrm{Pr}(\ulcorner G\urcorner) \to \mathrm{Pr}(\ulcorner
        \bot\urcorner)$.
\end{enumerate}
$\mathit{constr5}$ is the explicit primitive recursive function that
packages the three pieces into a single proof-tree transformer; the
chain $f\_\mathit{part}$, $g\_\mathit{part}$, $K\_\mathit{part}$,
$h\_\mathit{part}\_\mathit{skel}$ in Section~\ref{subsec:num-cor} above
is its syntactic decomposition.

So the connection to D2 is real but modest: D2 is one of three
ingredients, and its BRA realiser is just the proof-checker, not a
deep theorem.  Step 5 is \emph{not} L\"ob's theorem
($\Box(\Box A \to A) \to \Box A$); L\"ob is never used in the chain.

\paragraph{Modal-logic sidebar.}
For readers familiar with intuitionistic modal $S4$ in
Davies--Pfenning's style~\cite{daviespfenning2001}: the picture
above has a clean modal reading.  Their typed calculus has two
contexts ($\Gamma$ for ordinary value variables and $\Delta$ for
``code'' or ``boxed'' variables) and a type $\Box A$ of codes for
$A$, with $\mathbf{box}\,M : \Box A$ allowed only when $M$ depends
solely on $\Delta$-variables.  Under this reading: the $\mathit{cor}$
operation is the unit $A \to \Box A$; Theorem 12 is the action of
$\Box$ on a morphism (so $D_f$ is $\Box f$); Theorem 13 is
necessitation D1 for equations; Theorem 14 step 5 invokes D2 at $G$
plus the diagonal.  The notational gotcha above corresponds to the
fact that $\mathbf{box}\,(f\,x) : \Box B$ is well-typed but
$f\,(\mathbf{box}\,x)$ is a type error --- the modal type system
catches by inspection what BRA's uniform $\mathrm{Tree}$ type lets the
notation hide.  The modal reading is a useful specification, not a
shortcut: Davies--Pfenning's $\Box$ is metatheoretic and gives no
algorithmic content for $D_f$, $\mathit{constr5}$, or $\mathit{thmT}$.
G\"odel~II proper requires the provability predicate to be a primitive
recursive function definable inside BRA itself; the deep encoding is
intrinsic to the theorem.

\section{Methodology}

\subsection{Fast typecheck as a correctness oracle}

Every per-file solo typecheck runs in well under one second (with the
historical exception of \verb|BRA/Thm/ThmT.agda| at $\sim$19s; this is
the 16k-line dispatcher cascade and was already in place before the
G\"odel~II push).  The session followed a strict rule: \emph{slow
typecheck = mathematical mismatch, stop and re-architect}.  In
practice this prevented two specific traps:
\begin{enumerate}
  \item Trying to feed Agda an opaque normalisation it cannot perform
        (e.g.\ unfolding $\mathit{codeFormula}\;G$ at a parametric
        substituent position).  The fix is always: add a single
        targeted refl-witnessable lemma inside the right
        \verb|abstract| block (where the opaque definition unfolds),
        and use it as a black box outside.
  \item Forcing structural recursion through opaque trees
        ($\mathit{codedSubst} / \mathit{codedSubstT}$ on
        $\mathit{codeF}_1\;\mathit{thmT}$, etc.).  Same fix: prove the
        equation once inside the abstract block, expose only the
        type signature outside.
\end{enumerate}

\subsection{Architectural choice for $\mathit{constr5}$}

The original \verb|Th14Constr5Final.constr5| wraps $h\_\mathit{part}$
with two \emph{additional} encoded-mp layers
($h\_\mathit{part}\_\mathit{thxLayer}$, $h\_\mathit{part}$).  Under the
$\mathit{thmT}$ dispatcher these would consume both antecedents of
$h\_\mathit{part}\_\mathit{skel}$'s ex-falso shape, collapsing the value
of $\mathit{thmT}\;(h\_\mathit{part}\;x)$ to $\bot$ before it can serve
as antecedent to $\mathit{constr5}\_{\mathrm{inner}}$.  The fix is to
use $h\_\mathit{part}\_\mathit{skel}$ directly in
$\mathit{constr5}_{\mathrm{final}}$ (pure sb-chain, no internal mp's),
matching Guard's intent on page~17 and the comment in
\verb|Th14Constr5.agda|.  The two outer $\mathit{mp}$'s in
$\mathit{constr5}_{\mathrm{final}}$ then consume the two antecedents
exactly as Guard's chain prescribes.

\subsection{The Carneiro-lifted chain: simulating the deduction theorem in a Hilbert system}
\label{subsec:carneiro}

BRA's $\mathit{Deriv}$ is a Hilbert-style system: implication is a
connective, but there is no native introduction rule that, given a
derivation of $\psi$ depending on a hypothesis $\varphi$, returns
$\mathit{Deriv}\;(\varphi \mathrel{\mathit{imp}} \psi)$.  The propositional
axioms are Guard's
\[
  \mathit{axK} \;=\; P \mathrel{\mathit{imp}} (Q \mathrel{\mathit{imp}} P), \qquad
  \mathit{axS} \;=\; (P \mathrel{\mathit{imp}} (Q \mathrel{\mathit{imp}} R))
                       \mathrel{\mathit{imp}}
                       ((P \mathrel{\mathit{imp}} Q) \mathrel{\mathit{imp}} (P \mathrel{\mathit{imp}} R)),
\]
together with modus ponens.  $\mathit{axS}$ is Frege's $1879$ axiom.

In such a system the deduction theorem can be \emph{simulated} (a
technique attributed to Mario Carneiro in the metamath community):
replace every theorem $\psi_i$ in a derivation sequence by
$\varphi \mathrel{\mathit{imp}} \psi_i$.  If $\psi_k$ follows from $\psi_i$
and $\psi_j = \psi_i \mathrel{\mathit{imp}} \psi_k$ by modus ponens, then
$\varphi \mathrel{\mathit{imp}} \psi_k$ follows from
$\varphi \mathrel{\mathit{imp}} \psi_i$ and
$\varphi \mathrel{\mathit{imp}} (\psi_i \mathrel{\mathit{imp}} \psi_k)
 = \varphi \mathrel{\mathit{imp}} \psi_j$
via $\mathit{axS}$ and two applications of modus ponens.  Frege's
axiom is precisely what licenses the lifting; the form was apparently
understood in $1879$ and lost in Russell's and Hilbert's later
presentations.

The Theorem~14 chain uses this trick at the level of the running
hypothesis
\[
  \mathit{PrAtX}\;x \;:=\; (\mathit{thmT}\;x \;=\; \mathit{cor}\;G).
\]
Each step lemma $\mathit{step1\_l}, \dots, \mathit{step5\_l},
K\_\mathit{part}\_l, h\_\mathit{part}\_\mathit{skel}\_l$ is the
$\mathit{PrAtX}\;x \mathrel{\mathit{imp}} (\dots)$-lifted form of the
underlying step equation; the suffix ``\verb|_l|'' (\emph{lifted}) marks
this throughout the codebase.  Modus ponens between two adjacent step
lemmas is mediated by $\mathit{axS}$.  The reusable combinators sit in
\verb|BRA/Thm/EvalHelpers.agda| (header ``Lift helpers
(deduction-theorem combinators)''); the lifted Theorem~13 for singular
functors is in \verb|BRA/Th14Constr5.agda| (``Carneiro-lifted
Theorem~14 chain''); the per-step proofs live in
\verb|BRA/Th14Step3.agda|, \verb|Th14Step4Final.agda|,
\verb|Th14Step5.agda|, and \verb|Th14Step5HPart.agda|.  Theorem~12's
$\mathit{Rec}$~/~$\mathit{RecP}$ cases use the same idea at
\emph{dispatcher granularity} --- the per-tag dispatcher is what is
lifted, not each individual equation
(\verb|BRA/Th12RecP.agda|, \verb|BRA/TreeEqReflParam.agda|, with
helpers in \verb|BRA/Thm/EvalHelpers.agda|).

Without this simulation the Theorem~14 development would not close:
each step genuinely depends on $\mathit{PrAtX}\;x$, and a Hilbert-style
$\mathit{Deriv}$ offers no native means to discharge such a hypothesis.
The Carneiro lifting is therefore not a stylistic choice but a
structural component of the proof, on a par with the
$\mathit{cor}$/$\mathit{code}$ asymmetry of
Section~\ref{subsec:num-cor}.

\section{Statistics}

\subsection{File-by-file summary of the G\"odel~II push}

\begin{center}
\begin{tabular}{lrr}
\hline
\textbf{File} & \textbf{LoC} & \textbf{Solo typecheck} \\
\hline
\verb|BRA/Th14SubTLeaf.agda|       & 374 & 0.18\,s \\
\verb|BRA/Th14Step3Canon.agda|     & 582 & 0.60\,s \\
\verb|BRA/Th14Step3.agda|          & 633 & 0.81\,s \\
\verb|BRA/Th14Step4Final.agda|     & 213 & 0.84\,s \\
\verb|BRA/Th14Step5HPart.agda|     & 394 & 0.84\,s \\
\verb|BRA/Th14Step5.agda|          & 305 & 0.76\,s \\
\verb|BRA/GoedelIIFull.agda|       & ~50 & 0.71\,s \\
\hline
\end{tabular}
\end{center}

\subsection{Earlier infrastructure (already in place)}

\begin{itemize}
  \item \verb|BRA/Thm/ThmT.agda|: 16k LoC, 40-deep dispatcher cascade, 19\,s.
  \item \verb|BRA/Thm12.agda|: structural Theorem 12 with one Parts file per constructor, 0.65\,s.
  \item \verb|BRA/Thm13.agda|: $\sim$80 LoC, 0.1\,s.
  \item \verb|BRA/Thm11.agda|, \verb|Thm11Abstract|, \verb|Thm11Diagonal|: $\sim$300 LoC total.
  \item \verb|BRA/Thm14Abstract.agda|: $\sim$250 LoC, the $\mathit{closureToG}$ machinery.
  \item \verb|BRA/Thm14Constr5Final.agda|, \verb|Th14StepHPre|, \verb|Th14HUnfolds|, \verb|Th14GPartUnfolds|: $\sim$700 LoC of Pair-of-Comp2 unfold lemmas.
\end{itemize}

\subsection{Constraints maintained throughout}

\begin{itemize}
  \item \verb|--safe --without-K --exact-split| on every file.
  \item ASCII only.  No Unicode.
  \item No postulates.  No holes.  No \verb|with|-abstraction.  No dot patterns.
  \item No new $\mathit{Deriv}$ axioms beyond Guard's primitive list.
\end{itemize}

\section{Discussion}

The two G\"odel theorems in BRA, side by side:
\begin{itemize}
  \item \textbf{G\"odel I.} If the diagonal sentence $G$ is provable, then $\bot$ is provable.
  Equivalently (under the meta-assumption that BRA is consistent): $G$ is unprovable.
  \item \textbf{G\"odel II.} If the consistency formula $\mathit{Con}$ is provable, then $\bot$ is provable.
  Equivalently: BRA cannot prove its own consistency.
\end{itemize}

The intermediate functions $\mathit{constr5}_{\mathrm{final}}$ and the
parametric step lemmas $\mathit{step1\_l}$--$\mathit{step5\_l}$,
$K\_\mathit{part}\_l$, $h\_\mathit{part}\_\mathit{skel}\_l$ are
\emph{first-order} in the sense of recursive-arithmetic definability:
each is a closed combinator built from primitive $\mathit{Fun}_1$ /
$\mathit{Fun}_2$ symbols, and each $\mathit{Deriv}$ derivation is given
by an explicit equational proof in the BRA system.  No external
soundness theorem is invoked.

\subsection{Comparison with prior formalisations}

To our knowledge the only previous machine formalisation of G\"odel's
\emph{Second} Incompleteness Theorem is Paulson's Isabelle/HOL
development of hereditarily finite set theory~\cite{paulson2014}.
Other formalisations (O'Connor in Coq, Han in Lean, Carneiro's metamath
work) target G\"odel's \emph{First} theorem only.

Paulson's development uses HF set theory and full G\"odel numbering of
Hilbert-style proofs, with encoded substitution treated as an external
soundness theorem.  The present BRA development by contrast follows
Guard~\cite{guard1963}: the ``provability predicate''
$\mathit{thmT}$ is itself a primitive-recursive combinator built from
$\mathit{Rec} + \mathit{Fan} + \mathit{Lift} + \mathit{IfLf}$; encoded
substitution is performed by an explicit primitive recursive function
$\mathit{sub} : \mathit{Fun}_2$; and the entire chain (G\"odel I and
II) stays first-order, no Hilbert-proof G\"odel numbering, no external
soundness theorem.

\begin{thebibliography}{9}

\bibitem{guard1963}
J.~R.~Guard,
\newblock \emph{Lecture notes on recursive arithmetic},
\newblock Applied Mathematics Division, Argonne National Laboratory, 1963.

\bibitem{paulson2014}
L.~C.~Paulson,
\newblock \emph{A machine-assisted proof of G\"odel's incompleteness theorems for the theory of hereditarily finite sets},
\newblock Review of Symbolic Logic, vol.~7 no.~3, 2014, pp.~484--498.

\bibitem{daviespfenning2001}
R.~Davies and F.~Pfenning,
\newblock \emph{A modal analysis of staged computation},
\newblock Journal of the ACM, vol.~48 no.~3, May 2001, pp.~555--604.

\end{thebibliography}

\end{document}
